\section{A chiral index $\pm3$, given its inputs}\label{sec:three}

\paragraph{The question.}
The couplings of Section~\ref{sec:couplings} treat the triplet as a representation. A generation, in the physical
sense, is a \emph{count} of chiral fermions, which is an index. We ask when the grammar's triplet becomes such a count
and what must be assumed. The answer is conditional on stated inputs, and it is not an index of a thread. On any cusped
three-manifold or cover every twisted Euler characteristic vanishes (B1604; Boundary result~\ref{bd:noindex}), so no flat-module count on a single thread is a chiral index. The count lives on the fibre
family, as follows. The record is cited by arc~\cite{OriginAxiomRepo}.

\paragraph{The setting.}
The weave's surface is $M_{1,2}$, the moduli of a torus with two marked points, taken as the weave's surface by a reading made
within the program, of which the record verifies only the Euler characteristic (B1624). Its fibre loops, the point-pushes, are the inner automorphisms,
and the weave's moves lift to the matter by the lift that agrees with the ruled moves (B1641). That lift is the record's:
$\mathrm{conj}(a)$ and $\mathrm{conj}(b)$ act on the triplet by the parity signs and the commutator loop by $-1$, which is
$\Delta^4$ of the moves. The geometric lift, trivial on cohomology, differs from it by that central sign at $\Delta^4$.

\paragraph{The inputs.}
\begin{itemize}
\item \textbf{The triplet postulate} (Table~\ref{tab:inputs}): the left-handed matter lives in the triplet $T$ and the mass
  tensor couples it. A tagged working postulate, in the record the subspace named $\Lambda$ (renamed here to avoid the
  cosmological constant). What would discharge it: an even-dimensional object with a non-flat bundle whose index is the
  count.
\item \textbf{One cusp unit}: the cusp contributes $n$ units of the end datum, with $n=\pm1$ when one unit is taken.
  A tagged working postulate. What would discharge it: an end mechanism at the cusp.
\item \textbf{The extreme boundary domain}: the subspace $\Lambda_+$ of the critical coefficients at the cusp is taken to be
  $0$ or the whole module $L_6$. Other admissible domains give $\pm1$ (below). What would discharge it: a rule that selects
  the domain.
\end{itemize}
The setting, $M_{1,2}$ as the weave's surface with the lift that agrees with the ruled moves, is a reading verified only
in part, and we list it with the inputs.
None is derived in the record: the first is a postulate, the second a normalisation at an end the grammar does not fix,
the third a choice of domain. All are priced in the inputs table (Table~\ref{tab:inputs}). The ``$3$'' is the dimension
of the postulated triplet, read through the extreme domain; it is not a derived count of generations.

\begin{theorem}[The index $\pm3$; T6]\label{thm:index}
Given the triplet postulate and $n$ units at the cusp, on $M_{1,2}$ with the weave's lifts, the four-dimensional index
is $n(|\Lambda_+|-3)$, where $|\Lambda_+|$ is the dimension of the chosen subspace $\Lambda_+$ of the critical coefficients at the cusp. For
$\Lambda_+=0$ the index is $+3,0,-3$ for $n=-1,0,1$; for the two-dimensional piece $V_2$ of the six-dimensional module it
is $+1,0,-1$ and for $V_4$ it is $-1,0,+1$; and for the full module $L_6$ it is $-3,0,+3$. With one cusp unit, $n=\pm1$,
a non-zero (chiral) index is therefore $\pm3$ for $\Lambda_+=0$ or $\Lambda_+=L_6$.
\end{theorem}
\begin{proof}[Proof sketch]
The six-dimensional module $L_6$ is a representation of the metaplectic cover $\mathrm{Mp}_2(\Z)$ with exponents
$(1,1,3,5,7,7)/8$ and commutant of dimension $2$ (braid relation exact). It splits as $V_2\oplus V_4$, two non-isomorphic
irreducibles, so no three-dimensional constituent exists. The characters $\chi_1$ of the pieces and of $L_6$ vanish and
$\chi_{3/2}(\rho_T)=0=\chi_{1/2}(\bar\rho_T)$, so weight one is allowed, and the dimension formula for vector-valued modular forms
(as used in B1632) gives the fibre's contribution; the index is $n$ times it. The table is verified in both identifications of the module.
See B1641.
\end{proof}
\noindent\emph{Scope.} Verified for the module and the lifts named; the setting, $M_{1,2}$ as the weave's surface, is a
reading, and whether the lifts cohere over the whole of $\mathrm{Aut}^+(F_2)$ is open. The theorem is conditional on both inputs, and \emph{conditional} is the word: with $n=0$, or without the triplet
postulate, there is no count. B1641 also finds, by its sealed read-out, that the hand is \emph{not} the puncture
condition's orientation at the representation level: neither $\rho_T$ nor its conjugate is a sub or quotient of $L_6$. It
makes no claim about the hand (Section~\ref{sec:selection}).

\paragraph{The end radius, and what fixes it.}
Proposition~\ref{prop:norm} introduced the end radius $L=\log(1/\epsilon)$ of a disc removed at the puncture. B1656's
channel norm is the regularised norm with $2\pi L$ subtracted, and the first generation's lift $D/(2\pi L)$ depends on
it. In the flat metric a complete cusp has no boundary on which to impose a condition, and $L$ behaves as a free scale
(B1654). In the fibre's own complete hyperbolic metric, the uniformisation of the punctured torus, the puncture is a
hyperbolic cusp. A hyperbolic cusp has canonical neighbourhoods: the maximal embedded horoball, and the horocycle of
length $2$, which is embedded by Shimizu's lemma. Each fixes $L$ as a function of $\tau$ alone.

\begin{censuslaw}[the canonical cusp; B1667]\label{cl:cusp}
On a grid over the fundamental domain up to $\operatorname{Im}\tau=4$ ($54$ points), the maximal cusp gives
$L_{\max}\in[0.643,0.691]$ and the horocycle of length $2$ gives $L_2\in[2.29,2.78]$. At $\tau=e^{i\pi/3}$ both equal a
closed form from the periods of $\eta^4$: the punctured hexagonal torus is $\Hb/\SL(2,\Z)'$, freely generated by $LR$ and
$RL$. Under either canonical cut, the first generation's lift, computed with the exact channel norms outside the
horoball, is a function of $\tau$ alone. It is $0.03$ to $0.41$ in the bulk, has a simple zero at $\omega$, and falls like
$e^{-2.6\operatorname{Im}\tau}$ toward the cusp of moduli.
\end{censuslaw}

\noindent The end therefore supplies no hierarchy: a large $L$, which the bulk hierarchy $m_1/m_3\approx0.15/L$ would
need, is a cut far inside every canonical horoball, at a scale the fibre's geometry does not supply. The end datum is a
discrete choice between canonical horocycles, not a continuum, and a small first generation becomes a statement about
$\tau$. The computation is a census over the grid, informed by points seen before its seal. It does not say which
canonical horocycle, if either, is the physical end.

\paragraph{Does the symmetry fix the domain?}
If the domain must be preserved by the moves, the freedom drops from a continuum to the invariant subspaces of $L_6$. By
the proof above, $L_6=V_2\oplus V_4$ with $V_2$ and $V_4$ inequivalent and irreducible; the group the moves generate on
$L_6$ is the binary octahedral group $2O$, of order $48$ (B1606). So the invariant domains are $0,V_2,V_4,L_6$, and the
index is $\pm3$ or $\pm1$ (the record's kept indices $-3,-1,+1,+3$, B1606). The symmetry therefore
restricts the count but does not force three: an invariant domain of dimension three, which would leave only $-3,0,+3$,
does not exist, because $L_6$ has no three-dimensional constituent. The extreme domain remains a choice. Whether a
finite-distance model of the puncture, whose index is also $\dim\Lambda-3$, carries the same module is open.

\reading{The triplet is a representation until something counts it. What counts it is not in the grammar: one postulate
says where the matter lives, one unit says how many times the end is taken, and one choice says which boundary domain is
kept. A chiral three is thereby a question about the cusp's end, not a consequence of the moves.}
